\documentclass[11pt,a4paper]{article} %no openany perche' ci prendiamo spazio

\usepackage[utf8]{inputenc} %possono essere usati caratteri accentati

\usepackage[italian]{babel} %viene usato l'italiano

%evita che le immagini vengano spostate
\usepackage{float}
\usepackage{hyperref}
\hypersetup{
  colorlinks,
  citecolor=gray,
  filecolor=red,
  linkcolor=blue,
  urlcolor=blue
}

\usepackage{xcolor}
\usepackage{amsmath}
%immagini
\usepackage{graphicx}
\usepackage{caption}
%deprecated
\usepackage{subcaption}

%immagini in landscape
\usepackage{geometry}
\usepackage{pdflscape}

\usepackage{amsfonts}
\usepackage{listings}
\lstset{%
	commentstyle=\color{gray},
	frame=single,
	keepspaces=true,
	keywordstyle=\color{blue},
	numbers=left,
	numberstyle=\tiny\color{black},
	rulecolor=\color{black},
    basicstyle=\ttfamily
}

\usepackage{fancyhdr}


\title{\textbf{Relazione progetto Tecnologie Web}}
\author{ilQuartetto}
\date{5 febbraio 2016}


\newcommand{\makeFrontPage}{
  % Declare new goemetry for the title page only.
  \newgeometry{top=3.5cm}

  \begin{titlepage}
  \begin{center}

  \begin{center}
  \includegraphics[width=10cm]{res/img/logo.png}
  \end{center}

  \vspace{1cm}

  \begin{Huge}
  \textbf{Relazione progetto Tecnologie Web}
  \end{Huge}

  \textbf{\emph{Gruppo ilQuartetto, Progetto Passione Karaoke}}

  \vspace{11pt}

  \bgroup
  \def\arraystretch{1.3}
  \begin{tabular}{ r|l }
    \multicolumn{2}{c}{\textbf{Informazioni sul gruppo}} \\
    \hline
    1070162 & Davide Polonio \\
    1070541 & Fabio Massignan \\
    1070258 & Andrea Mantovani \\
    1073397 & Luca Bianco \\
  \end{tabular}
  \egroup

  \vspace{22pt}

  \end{center}
  \end{titlepage}

  % Ends the declared geometry for the titlepage
  \restoregeometry
}


\fancypagestyle{plain}{
	% cancella tutti i campi di intestazione e piè di pagina
	\fancyhf{}

	\lhead{
		\parbox[b]{10cm}{
			\emph{ilQuartetto} \vspace{0pt} \\
			\emph{Progetto PassioneKaraoke} \vspace{7pt}
		}
	}
	\chead{}
	\rhead{
		\slshape \leftmark
	}

	\rfoot{Pagina \thepage{}}

	% Visualizza una linea orizzontale in cima e in fondo alla pagina
	\renewcommand{\headrulewidth}{0.3pt}
	\renewcommand{\footrulewidth}{0.3pt}

}

\setlength{\headheight}{30pt}
\pagestyle{plain}


\begin{document}

\makeFrontPage

\newpage

\tableofcontents

%da qui in poi ci andranno gli include dei vari capitoli
\section{Introduzione}

\subsection{Abstract}
L`obiettivo che questo progetto si pone \`e l'implementazione di un sito web dedicato alla consultazione di testi di canzoni (in inglese \textit{lyrics}), e di articoli inerenti al mondo musicale. 
Dalla pagina principale, il visitatore pu\`o accedere alla sezione relativa agli Artisti, a quella relativa agli Articoli o alla pagina del log in.
La sezione Artisti, nello specifico, contiene tutti i cantanti, e per ogni cantante i relativi album. Selezionando poi l'album di interesse, si potr\`a accedere all'elenco delle canzoni in esso contenute e, per ogni canzone, saranno disponibili il testo e il video del karaoke.
Per quando riguarda gli Articoli, questi si potranno consultare nella relativa sezione, dove compariranno in ordine cronologico a partire dal pi\`u recente.
Infine effettuando il log in, l'Utente registrato accede alla sua pagina personale in cui ritrova le sue canzoni preferite. L'Utente registrato ha infatti la possibilit\`a di esprimere un giudizio positivo o negativo per ogni canzone; e le canzoni per cui esprime un giudizio positivo vengono aggiunte direttamente nella pagina personale, dove possono essere comodamente consultate dall'Utente anche in momenti successivi.
L'Amministratore è un tipo particolare di Utente che pu\`o ampliare il catalogo aggiungendo nuovi Artisti, con i relativi album e canzoni, oppure nuovi album o nuove canzoni in Artisti e Album gi\`a esistenti. Può inoltre aggiornare gli utenti sulle novità del mondo della musica tramite la creazione di news (articoli).


\newpage %interruzione di pagina


\subsection{A chi è rivolto il sito}
Il sito sar\`a consultato da utenti molto vari, in quanto non tratta una materia di nicchia ma \`e inerente a un mondo di interesse pressoch\`e universale com'\`e quello della musica. Questo grande bacino di utenti comporta, di conseguenza, la necessit\`a di porre particolare attenzione all'accessibilit\`a generale nella fase di implementazione. Proprio con l'obiettivo di migliorare l`accessibilit\`a, per esempio, i contenuti saranno inseriti da un \textbf{Admin}, che si occuperà di scrivere articoli in formato \textit{HTML}, rispettandone tutte le convenzioni e inserendo gli opportuni tag per rendere la lettura dei contenuti fruibile anche agli utenti che hanno problemi visivi.


\subsection{Credenziali d'accesso al sito}
Per accedere al sito sono state creade delle credenziali di default.
\begin{itemize}

    \item Livello admin:
    \begin{itemize}

        \item Utente: \textit{admin}
        \item Password: \textit{password}

    \end{itemize}

    \item Livello utente:
    \begin{itemize}

        \item Utente: \textit{utente}
        \item Password: \textit{utente}

    \end{itemize}

\end{itemize}

\newpage %interruzione di pagina per staccare i contenuti
\section{Progettazione}

Di seguito vengono viste e spiegate alcune scelte di design effettuate durante la progettazione del sito.

\subsection{Accessibilità}
%TODO
Il sito è stato utilizzato sviluppando solamente codice XHTML 1.0 Strict, per assicurare la massima compatibilità. Per i fogli di stile è invece stato deciso di applicare CSS 3.0, per potersi avvantaggiare di nuove tecnologie grafiche e ridurre al minimo l'utilizzo di JavaScript, che invece è stato applicato per i controlli sul form di login e sulla gestione del contenuto del sito, avendo comunque una buona degradazione in caso questa tecnologia non sia disponibile nel dispositivo in uso.

\subsection{Usabilità}
%TODO

\subsection{Perl}
%TODO

\subsection{HTML}
%TODO

\subsection{CSS}
%TODO

\subsection{Javascript}
%TODO

\newpage
\section{Appendice}

\subsection{Organizzazione del gruppo}

Il lavoro nel progetto in esame è stato così distribuito:

\begin{itemize}

    \item Fabio Massignan
    \begin{itemize}

        \item Progettazione e realizzazione dei file XML e XMLSchema
        \item Stesura HTML
        \item Stesura relazione
        \item Recupero del contenuto, imamgini, icone, testi
    \end{itemize}

    \item Davide Polonio
    \begin{itemize}

        \item Stesura HTML
        \item Progettazione e stesura relazione
        \item Progettazione e realizzazione JavaScript

    \end{itemize}

    \item Andrea Mantovani
    \begin{itemize}

        \item Progettazione e realizzazione Perl
        \item Progettazione e realizzazione JavaScript
        \item Stesura relazione
        \item Stesura HTML

    \end{itemize}

    \item Luca Bianco
    \begin{itemize}

        \item Progettazione e realizzazione CSS
        \item Stesura relazione

    \end{itemize}


\end{itemize}



\end{document}
